\section{Results}

\subsection{Replication summary (20 seeds)}
\begin{table}[h]
\centering
\begin{tabular}{l c}
\toprule
Metric & Value \\
\midrule
Emergent count & $6.7 \pm 3.1$ (min 5, max 14) \\
Earlier-share & $1.000$ (std $0.000$) \\
Shuffle z-score & $2.96 \pm 0.86$ \\
Runs with $z>2$ & $20/20$ \\
\bottomrule
\end{tabular}
\caption{20-seed replication summary for temporal scaffolding statistics.}
\label{tab:replication20}
\end{table}

Table~\ref{tab:replication20} shows complete seed-level success under the shuffle null criterion ($z>2$), with zero variance in earlier-share across the 20-run panel.

\subsection{Deep-run characterization}
Deep run statistics: earlier-share $=1.000$, shuffle-null $z\approx7.20$, degree-preserving-null $z\approx96$.

Graph structure: 186 nodes, 17{,}205 edges, and 850 backbone edges.
Emergent scaffold: 104 nodes, 225 edges, earlier share $\approx0.791$, shuffle $z\approx6.5$.

Concentration metrics indicate heavy access inequality: Gini(access) $=0.943$ and Top10 share $=0.976$.

Two-timescale mixture fit yields short mode $\approx1190\,$s and long mode $\approx98{,}313\,$s, with $\Delta$BIC strongly favoring two components.

\begin{figure}[h]
  \centering
  \includegraphics[width=0.95\linewidth]{figures/fig_scaffold_directed.png}
  \caption{Emergent-only directed scaffold, retaining only newer$\rightarrow$older edges.}
  \label{fig:scaffold-directed}
\end{figure}

\begin{figure}[h]
  \centering
  \includegraphics[width=0.95\linewidth]{figures/fig_null_distributions.png}
  \caption{Null-model earlier-share distributions (shuffle and degree-preserving) with observed statistic marked.}
  \label{fig:null-distributions}
\end{figure}

\begin{figure}[h]
  \centering
  \includegraphics[width=0.8\linewidth]{figures/fig_age_gap_mixture.png}
  \caption{Histogram of $\log(\Delta t)$ with one-component and two-component mixture overlays, annotated by BIC contrast.}
  \label{fig:age-mixture}
\end{figure}

\begin{figure}[h]
  \centering
  \includegraphics[width=0.95\linewidth]{figures/fig_access_concentration.png}
  \caption{Access concentration diagnostics: Lorenz curve and rank-frequency plot.}
  \label{fig:access-concentration}
\end{figure}

\subsection{Robustness checks}
We ran two sensitivity checks to evaluate whether the temporal scaffolding signal is a byproduct of extraction choices.
First, backbone extraction was varied across $K\in\{4,6,8\}$, preserving the directional bias trend across settings rather than producing a single-$K$ artifact.
Second, we excluded short-mode edges (assigned to the fast timescale in the two-component fit) and re-estimated earlier-share against both null families; the backward-linking bias remained directionally consistent, though reduced sample size widened uncertainty bands.
These checks support robustness of the qualitative finding while keeping quantitative interpretation conservative.
